\section{Deletion operators and injectivity of the leading projection}
\label{sec:deletion}

Our goal is to show that setting $A=0$ loses no leading part
$\chi(A,C,B)$ of a reduced rational solution.  We use deletion
operators to test the four-term relation of
Proposition~\ref{prop:leading-pentagon}.  These tests first show that
a leading part killed by $A=0$ can involve only $A,B$; the special
identity then excludes this remaining kernel.

\subsection{The integral deletion representation}
The operators below specialize the integral algebraic differential
operators of Hirose and Sato \cite[Section~2]{HiroseSato}.
Their logarithmic-form and cyclotomic braid interpretation follows
Hirose's construction \cite[Section~4.3]{Hirose}.
Here we construct an integral representation of the signed-hyperplane
algebra of Section~\ref{sec:degeneration}.  After reduction modulo two it
will turn the leading pentagon into individual coefficient equations.
Throughout the integral construction, hyperplanes and their incidences
are taken over characteristic zero.

Following the ratio alphabet in \cite[Section~6]{Hirose}, put
\begin{equation}\label{eq:deletion-letters}
 e=1,\qquad
 p=\frac{z_2-z_3}{z_2+z_3},\qquad
 w=\frac{z_2-z_1}{z_2+z_1}.
\end{equation}
For an integer $N\geq0$, let $E_N$ be the free abelian group with basis the
words of length at most $N$ in the alphabet $\{0,e,p\}$, including the
empty word.  Every word has fixed left and right endpoints $0,w$; these
endpoints are never deleted.  If $P,Q\in\{0,e,p,w\}$ are distinct, and
$H\in\mathcal A_3$, write
\[
 \nu_H(P,Q)=\operatorname{ord}_H(P-Q).
\]
For equal letters we set $\nu_H(P,P)=0$, rather than taking the order of
the zero function.  Thus the convention is part of the definition of
the operators.

Define an endomorphism $T_H$ of $E_N$ by
\begin{equation}\label{eq:deletion-operator}
 T_H(P_1\cdots P_k)
 =\sum_{j=1}^{k}
 \bigl(\nu_H(P_j,P_{j+1})-\nu_H(P_{j-1},P_j)\bigr)
 P_1\cdots\widehat{P_j}\cdots P_k,
 \qquad P_0=0,\quad P_{k+1}=w.
\end{equation}
In particular $T_H$ kills the empty word.  Every $T_H$ is a finite
integer matrix.  The specialization from \cite[Section~2]{HiroseSato}
uses the field $\Q(z_1,z_2,z_3)$, the valuation $\operatorname{ord}_H$,
and endpoints $0,w$, with the same zero-difference convention.

\begin{proposition}\label{prop:deletion-representation}
The assignment $t_H\mapsto T_H$ defines a representation of
$\mathfrak h_{\Z}$ on $E_N$.  Equivalently, for every codimension-two
flat $X$ and every hyperplane $H_0\supset X$,
\begin{equation}\label{eq:deletion-flat}
 \left[T_{H_0},\sum_{H\supset X}T_H\right]=0.
\end{equation}
It therefore extends to the ordinary enveloping algebra, over $\Z$ and
after reduction over $\F$.
\end{proposition}

\begin{proof}
For distinct letters put $\omega_{PQ}=d\log(P-Q)$, and put
$\omega_{PP}=0$.  Signs of differences have no effect on these forms.
The nonconstant differences, up to sign, factor as
\begin{align}
 p&=\frac{z_2-z_3}{z_2+z_3},&
 w&=\frac{z_2-z_1}{z_2+z_1},\notag\\
 1-p&=\frac{2z_3}{z_2+z_3},&
 1-w&=\frac{2z_1}{z_2+z_1},\label{eq:deletion-factorization}\\
 p-w&=\frac{2z_2(z_1-z_3)}{(z_2+z_3)(z_2+z_1)}.&&\notag
\end{align}
The remaining nonzero difference $1-0$ is constant.  Consequently every
$\omega_{PQ}$ is an integer linear combination of $d\log\ell_H$, where
$\ell_H$ is a linear defining equation of $H$.

Consider the matrix-valued rational one-form
\[
 \Omega=\sum_{H\in\mathcal A_3}T_H\,d\log\ell_H.
\]
The coefficient of a deletion at $P_j$ is exactly
$\omega_{P_jP_{j+1}}-\omega_{P_{j-1}P_j}$.  We claim that
$\Omega\wedge\Omega=0$.  For two nonadjacent deleted positions, neither
deletion changes the neighbors of the other, so the contributions of
the two orders cancel by antisymmetry of the wedge product.  For two
adjacent positions, denote the four successive letters by $L,P,Q,R$.
The sum of the contributions, with the rightmost matrix acting first,
is
\begin{align}
 &(\omega_{QR}-\omega_{LQ})\wedge
       (\omega_{PQ}-\omega_{LP})\notag\\
 &\hspace{15mm}+
 (\omega_{PR}-\omega_{LP})\wedge
       (\omega_{QR}-\omega_{PQ}).\label{eq:adjacent-deletions}
\end{align}
Expanding this expression cancels the cross terms involving
$\omega_{LP}$ and $\omega_{QR}$.  The remaining terms are the two
logarithmic triangle identities for $(L,P,Q)$ and $(P,Q,R)$.  The
triangle identity is
\[
 \omega_{IJ}\wedge\omega_{JK}
 +\omega_{JK}\wedge\omega_{KI}
 +\omega_{KI}\wedge\omega_{IJ}=0.
\]
For three distinct letters it follows by differentiating the relation
$(I-J)+(J-K)+(K-I)=0$ and putting the terms over a common denominator.
If exactly two letters agree, one form is zero and the other two are
equal; if all three agree, all three forms are zero.  Thus it also holds
in every repeated-letter case.  This proves that
\eqref{eq:adjacent-deletions} vanishes, including when other pairs among
$L,P,Q,R$ agree.  It proves the claimed matrix identity.

Take its residue first along a fixed $H_0$.  Terms involving no
$d\log\ell_{H_0}$ have no pole at a generic point of $H_0$, and the
result is the rational one-form identity on $H_0$
\[
 \sum_{H\ne H_0}[T_{H_0},T_H]\,
 d\log(\ell_H|_{H_0})=0.
\]
Now take the residue on $H_0$ along a codimension-one subspace
$X\subset H_0$ which is a flat of the arrangement.  For $H\ne H_0$,
the restricted nonzero linear form $\ell_H|_{H_0}$ has order one along
$X$ precisely when $H\supset X$, and has order zero otherwise.  This
gives \eqref{eq:deletion-flat}; including $H_0$ in the sum adds a zero
commutator.

This residue calculation is performed over $\Q$.  Distinct hyperplanes
may restrict to multiples of the same equation on $H_0$; each matrix
then contributes separately with coefficient one.  In particular,
restrictions proportional to $z$, $-z$, or $2z$ all have residue one.
There is no division by the number of coincident restrictions, nor by
the constant $2$.  Finally every matrix in \eqref{eq:deletion-flat} has
integer entries.  Its vanishing over $\Q$ is therefore its vanishing
over $\Z$, which permits reduction modulo two without changing the
characteristic-zero incidence relations.
\end{proof}

For a word $s$ of length $\ell$, and an enveloping-algebra element $q$
of length $\ell$, choose $N\geq\ell$ and define its scalar deletion test by
\[
 \langle q,s\rangle=[\varnothing]\,T(q)s,
\]
where $[\varnothing]$ extracts the coefficient of the empty word and
$T(q_1\cdots q_\ell)=T(q_1)\cdots T(q_\ell)$.  Thus the rightmost
factor acts first.  This convention will make successive deletion at
the right endpoint read an associative coefficient in its usual
left-to-right order.

\begin{lemma}\label{lem:edge-table}
In the integral generator basis
$a_i=t_{z_i=0}$, $u_{ij}=t_{z_i-z_j=0}$,
$v_{ij}=t_{z_i-z_j=0}+t_{z_i+z_j=0}$, the complete nonzero edge weights
are the following linear functionals:
\begin{equation}\label{eq:deletion-edge-table}
\begin{array}{c|c|c}
\text{edge}&\text{$a,v$ part}&\text{$u$ part}\\ \hline
0p&0&u_{23}^{*}\\
0w&0&u_{12}^{*}\\
0e&0&0\\
ep&a_3^{*}-v_{23}^{*}&0\\
ew&a_1^{*}-v_{12}^{*}&0\\
pw&a_2^{*}+v_{13}^{*}-v_{12}^{*}-v_{23}^{*}&u_{13}^{*}
\end{array}
\end{equation}
Equal-letter edges have weight zero.  After reduction modulo two, a
word containing $r$ zeros has zero scalar test on every product with
fewer than $r$ factors of type $u$.  For a product with exactly $r$ such
factors, every $u$-factor must delete a zero in any nonzero contribution.
\end{lemma}

\begin{proof}
If a difference has orders $\lambda_-$ and $\lambda_+$ along
$z_i-z_j=0$ and $z_i+z_j=0$, its evaluations on $u_{ij}$ and $v_{ij}$
are respectively $\lambda_-$ and $\lambda_-+\lambda_+$.  Applying this
to \eqref{eq:deletion-factorization} gives every entry of
\eqref{eq:deletion-edge-table}.  For example $p$ has orders $1,-1$,
so its $v_{23}$ weight is zero already integrally; the factor $2$ in
$p-w$ contributes no hyperplane valuation.

Every $a$- or $v$-weight on an edge touching $0$ is zero.  Thus an
$a$- or $v$-operation cannot delete a zero, regardless of preceding
deletions, because all remaining letters and endpoints still belong to
$\{0,e,p,w\}$.  Each $u$-operation removes at most one zero.  Emptying
$r$ zeros consequently requires at least $r$ such operations.  If there
are exactly $r$, using even one to remove a nonzero letter makes it
impossible to remove the remaining zeros.  The same conclusion holds
for linear combinations within the indicated slots.
\end{proof}

\subsection{Detecting mixed terms}
We now work over $\F$.  Let $\bar\psi\in\overline L_n$ have leading
$B$-count $r$ and leading part $\chi(A,C,B)$, with $n=2m+r>2$.
Every word of $\chi$ has $m$ letters from $\{A,C\}$ and $r$ letters
$B$.  Proposition~\ref{prop:leading-pentagon}, tested on a word with
$m$ nonzero internal letters and $r$ zeros, is an exact scalar identity:
the lower-count remainder is killed by Lemma~\ref{lem:edge-table}.
We use this exact identity to eliminate every component involving
both $A$ and $C$ from a possible kernel of the projection $A=0$.

\begin{lemma}\label{lem:deletion-kernel}
If $\chi(0,C,B)=0$, then $\chi$ belongs to the ordinary free Lie
algebra on $A,B$; in particular it has the form $P(A,B)$ with exactly
$m$ letters $A$ and $r$ letters $B$.
\end{lemma}

\begin{proof}
The associative multihomogeneous components of an ordinary Lie
polynomial are again ordinary Lie polynomials.  Since the component
with no $A$ vanishes, suppose that $h>0$ is the least $A$-count of a
nonzero component $\chi_h$.  We will rule out $h<m$.

Fix any internal word $s$ with $h$ letters $e$, $m-h$ letters $p$, and
$r$ zeros.  In each summand of \eqref{eq:leading-pentagon}, an
$A$-, $C$-, or $B$-slot means the corresponding first, second, or third
argument of $\chi$.  Every $B$-slot must delete a zero by
Lemma~\ref{lem:edge-table}.  The other relevant actions follow directly
from the edge table; the four rows are in the order of that equation:
\begin{equation}\label{eq:deletion-slot-actions}
\begin{array}{c|c|c|c}
 &A\text{-slot}&C\text{-slot on nonzero letters}&B\text{-slot on zeros}\\ \hline
1&\text{terminal }e&\text{terminal }e&\text{terminal }0\\
2&\text{terminal }e&
 \substack{p\text{ adjacent to }e\text{ or }w;\\
             e\text{ adjacent to }p}&0\text{ adjacent to }p\\
3&0&\text{unneeded}&\text{unneeded}\\
4&\text{terminal }e&\text{terminal }e\text{ or }p&\text{terminal }0
\end{array}
\end{equation}
Here ``terminal'' means the last remaining internal letter, adjacent to
$w$.  Each indicated boundary contributes weight one, so two such
boundaries contribute zero.  For instance, the third $A$-slot is
$a_1+a_2+v_{12}$: its weights on $ep,ew,pw$ are respectively
$0,1+1,1+1$, and its weights on edges touching zero vanish.  Its deletion
operator is identically zero.  In the second row the $C$-slot is
$v_{23}$ and the $B$-slot is $u_{23}$.

A nonzero contribution from an $A$-count-$t$ component must use each
$A$-slot to delete an $e$, so $t\leq h$.  Minimality gives $t\geq h$.
Hence only $\chi_h$ can contribute, and all $h$ letters $e$ must be
deleted by its $A$-slots.  In particular, the possible $C$-deletions of
$e$ in the table cannot occur in a surviving contribution.  This also
excludes contributions from every larger $A$-count.

In the fourth summand every deletion is therefore at the right
endpoint, and a slot word has at most one complete deletion sequence.
It succeeds precisely when it is the word $s'$ obtained by
$e\mapsto A$, $p\mapsto C$, $0\mapsto B$, in which case its weight is
one.  Thus the fourth summand evaluates to the coefficient $[s']\chi_h$.
The first summand cannot delete any $p$, so it vanishes because
$m-h>0$.  The third summand vanishes because $h>0$ and its $A$-operator
is zero.

Consider the second summand.  Suppose an $e$ is followed by a nonempty
maximal block consisting of $p$'s and zeros, ending at the next $e$ or
at $w$.  Its left $e$ cannot disappear while any of this block remains:
all $e$'s must be removed by $A$-slots, which remove only a terminal
$e$.  If the right boundary is an $e$ and that $e$ disappears before
the block does, it must already be terminal, so the new right boundary
is $w$.  Thus a surviving nonempty part of the block always has
boundaries $e,e$ or $e,w$.  The final surviving letter of this block
cannot be removed.  If that letter is $p$, its $C$-weight is $1+1=0$;
if it is $0$, its $B$-weight
is $0+0=0$.  Consequently the second summand has no complete deletion
sequence whenever some $e$ occurs before a non-$e$.
Figure~\ref{fig:deletion-block} records the two possible boundaries and
the last-letter obstruction.

\begin{figure}[tbp]
  \centering
  \begin{tikzpicture}[
    x=1cm,y=1cm,
    letter/.style={minimum width=.48cm,minimum height=.48cm,inner sep=1pt},
    boundary/.style={letter,draw,line width=.6pt,rounded corners=1pt},
    block/.style={draw=black!55,fill=black!4,rounded corners=2pt,
      minimum width=3.25cm,minimum height=.68cm,inner sep=3pt},
    move/.style={-{Stealth[length=2mm]},line width=.65pt},
    note/.style={font=\small,align=center,inner sep=2pt}
  ]
    \node[note,anchor=west] at (-5.7,1.15)
      {A nonempty block and its boundaries};
    \node[boundary] (left) at (-4.7,0) {$e$};
    \node[block] (word) at (-2.45,0) {$p\;0\;\cdots\;p\;0$};
    \node[boundary] (right) at (-.2,0) {$e$};
    \node[letter] (end) at (1.0,0) {$w$};
    \draw (left) -- (word) -- (right) -- (end);
    \draw[decorate,decoration={brace,amplitude=4pt,mirror}]
      (-4.0,-.48) -- (-.9,-.48)
      node[note,midway,below=5pt] {letters in $\{p,0\}$};
    \node[note,align=left,anchor=west] at (1.75,0)
      {If the right $e$ is terminal,\\it may be removed.};

    \draw[move] (-2.45,-1.2) -- (-2.45,-1.85);
    \node[boundary] (leftnext) at (-4.7,-2.35) {$e$};
    \node[block] (wordnext) at (-2.45,-2.35) {$p\;0\;\cdots\;p\;0$};
    \node[boundary] (rightnext) at (-.2,-2.35) {$w$};
    \draw (leftnext) -- (wordnext) -- (rightnext);
    \node[note,align=left,anchor=west] at (1.75,-2.35)
      {The left $e$ is still not terminal:\\it cannot be removed first.};

    \draw[move] (-2.45,-2.95) -- (-2.45,-3.6)
      node[note,midway,right=4pt] {any further deletions};
    \node[note,anchor=west] at (-5.7,-4.0)
      {The final surviving block letter, with $q\in\{e,w\}$};

    \node[boundary] (ep) at (-4.7,-4.8) {$e$};
    \node[letter] (p) at (-3.5,-4.8) {$p$};
    \node[boundary] (qp) at (-2.3,-4.8) {$q$};
    \draw (ep) -- node[above,font=\small] {$1$} (p)
      -- node[above,font=\small] {$1$} (qp);
    \node[note] at (-3.5,-5.5) {$C$-weight: $1+1=0$};

    \node[boundary] (ez) at (.6,-4.8) {$e$};
    \node[letter] (zero) at (1.8,-4.8) {$0$};
    \node[boundary] (qz) at (3.0,-4.8) {$q$};
    \draw (ez) -- node[above,font=\small] {$0$} (zero)
      -- node[above,font=\small] {$0$} (qz);
    \node[note] at (1.8,-5.5) {$B$-weight: $0+0=0$};
  \end{tikzpicture}
  \caption{The obstruction in the second summand of the leading pentagon,
  after the minimal $A$-count forces every $e$ to be removed by an
  $A$-operation. Such an operation removes only a terminal $e$.
  While the displayed $p,0$ block is nonempty, its left boundary remains
  $e$; removing a terminal right boundary changes $e$ to $w$.
  With either right boundary, the final surviving $p$ or $0$ has deletion weight
  zero in $\mathbb F_2$. Thus no nonzero deletion sequence empties the block.}
  \label{fig:deletion-block}
\end{figure}


The scalar leading pentagon now gives $[s']\chi_h=0$ for every such
$s$.  The only remaining words have the form $v e^h$, with $v$ a word
of $p$'s and zeros.  Because $m-h>0$, the word $v$ is nonempty.  Let
$\operatorname{rev}$ denote reversal of associative words.  For every
ordinary homogeneous Lie polynomial $q$ of associative length $d$,
\[
 \operatorname{rev}(q)=(-1)^{d-1}q.
\]
Indeed reversal sends $[q_1,q_2]$ to
$-[\operatorname{rev}(q_1),\operatorname{rev}(q_2)]$, proving the
formula by induction on Lie monomials.  In characteristic two it fixes
$\chi_h$.  The coefficient of $v e^h$ therefore equals the coefficient
of $e^h\operatorname{rev}(v)$, which has already been shown to vanish.
Every coefficient of $\chi_h$ is zero, a contradiction.

Thus a least nonzero $A$-count can only be $m$.  All $A$-counts are at
most $m$, so only the component using $A,B$ remains.  It is an ordinary
Lie polynomial as asserted.
\end{proof}

\subsection{Excluding the residual kernel}
Deletion has reduced the possible kernel to polynomials $P(A,B)$.
To exclude them, we now use the special identity rather than further
deletion tests.  Its lowest-count part compares coefficients obtained
by exchanging an $A$ and a $B$; antisymmetry will handle the remaining
case with only one $A$.

\begin{proposition}\label{prop:projection-injective}
For $n>2$ and every $r\geq0$, the leading projection
\[
 \operatorname{gr}^{r}\overline L_n\longrightarrow
 \Lie_{\F}\langle C,B\rangle,\qquad
 \chi\longmapsto\chi(0,C,B)
\]
is injective.  Its source is zero unless $n=2m+r$ for an integer
$m\geq0$; in that case its image has $C$-count $m$ and $B$-count $r$.
\end{proposition}

\begin{proof}
The map is well defined because changing a representative by an
element of $F^{r+1}$ does not change its count-$r$ component.  The
parity assertion and the multidegree of its image follow from
$\wt(A)=\wt(C)=2$ and $\wt(B)=1$.  If $m=0$, the leading part is an ordinary Lie polynomial
in $B$ alone and is zero in the weights at issue.  We may therefore
assume $m>0$.

Suppose that $\chi(0,C,B)=0$.  By Lemma~\ref{lem:deletion-kernel},
$\chi=P(A,B)$.  If $r=0$, this is an ordinary Lie polynomial in $A$
alone, and hence zero except possibly in weight two.  Thus a nonzero
kernel would require $r\geq1$.

Apply the reduction of Lemma~\ref{lem:special} and antisymmetry with
$z=x+y$.  They give
\begin{equation}\label{eq:deletion-special}
 [x,\bar\psi(x,y)+\bar\psi(z,y)]+[y,\bar\psi(z,y)]=0.
\end{equation}
Indeed the special identity with $a=x$, $b=y$, $c=z$ becomes
$[x,\bar\psi(y,x)]+[z,\bar\psi(y,z)]=0$, and antisymmetry in
characteristic two interchanges each pair of arguments.

Under $x\mapsto x+y$, the three associative elements become
\[
 A\mapsto A+C+B^2,\qquad C\mapsto C,\qquad B\mapsto B.
\]
Moreover, on their free associative algebra,
\[
 [x,A]=0,\qquad [x,B]=C,\qquad [x,C]=[A,B].
\]
The part of $\ad_x$ that lowers $B$-count by one is consequently the
derivation
\[
 D(A)=D(C)=0,\qquad D(B)=C;
\]
its other part raises $B$-count by one.  The substitution
$A\mapsto A+C+B^2$ never lowers that count.  Therefore components of
$\bar\psi$ with count greater than $r$ cannot contribute to count
$r-1$ in \eqref{eq:deletion-special}, and its outside bracket with $y$
has count at least $r+1$.  Within the leading part, using any $B^2$
in the substitution also raises the count by two.  The count-$r-1$
component of \eqref{eq:deletion-special} is exactly
\begin{equation}\label{eq:residual-kernel-identity}
 D\bigl(P(A+C,B)+P(A,B)\bigr)=0.
\end{equation}

Write $c(v)$ for the coefficient of a word $v$ with $m$ letters $A$
and $r$ letters $B$ in $P$.  Fix any target word with $m-1$ letters
$A$, $r-1$ letters $B$, and two letters $C$.  Its coefficient on the
left of \eqref{eq:residual-kernel-identity} is the sum of exactly two
coefficients of $P$: fill the two marked $C$-positions with $(A,B)$
or with $(B,A)$.  In each case the original $A$ is replaced by $C$
in the substitution and the original $B$ by $C$ under $D$; these are
the only possibilities.  The term $D(P(A,B))$ has only one $C$ and
contributes nothing here.  Thus the coefficients of $P$ are invariant
under exchanging any chosen $A$-position and $B$-position.  These
exchanges connect all words with the specified multiplicities, so
\begin{equation}\label{eq:constant-residual-coefficients}
 c(v)=\lambda\quad\text{for every such word }v
\end{equation}
for one $\lambda\in\F$.

If $m\geq2$, fix a target word with $m-2$ letters $A$, $r-1$ letters
$B$, and three letters $C$.  Exactly one of its three $C$-positions
must have been the $B$ differentiated by $D$; the other two must have
been $A$'s replaced in the substitution.  There are exactly three
choices.  Its coefficient is therefore $3\lambda=\lambda$, which
\eqref{eq:residual-kernel-identity} forces to vanish.  Hence $P=0$.

It remains to consider $m=1$.  The polynomial $P$ has original
$(x,y)$-bidegree $(2,r)$.  An $A,C,B$ word with counts $(a,c,b)$ has
original bidegree $(2a+c,c+b)$.  Solving for bidegree $(2,r)$ gives
only
\[
 (a,c,b)=(1,0,r)\quad\text{or}\quad(0,2,r-2).
\]
The latter possibility is excluded by the assumption
$\bar\psi\in F^r$, and the former is precisely $P$.  Thus if $P\ne0$,
the entire bidegree-$(2,r)$ component of $\bar\psi$ is nonzero.
Symmetry under $x\leftrightarrow y$ forces a nonzero component of
bidegree $(r,2)$.  Its $B$-count is at most its $y$-count, namely two,
so $F^r$ forces $r\leq2$.

For $r=2$, the ordinary free Lie component with one $A$ and two $B$'s
is spanned by
\[
 [B,[A,B]]=ABB+BBA.
\]
To see the spanning assertion directly, every three-letter bracketing
with two identical $B$'s either has the inner bracket $[B,B]=0$, or
is $[B,[A,B]]$ up to sign and skew-symmetry.  The coefficients on
$ABB,BAB,BBA$ are therefore $(\mu,0,\mu)$, which are all equal as
required by \eqref{eq:constant-residual-coefficients} only when
$\mu=0$.

For $r=1$, the forced swapped bidegree is $(1,2)$.  The count equations
$2a+c=1$, $c+b=2$ have the unique nonnegative solution
$(a,c,b)=(0,1,1)$.  Thus this nonzero component is a multiple of
$[C,B]$, lies in count one, and survives the projection $A=0$.
No higher-filtration component has this bidegree, so it cannot cancel
this projection.  This contradicts $\chi(0,C,B)=0$.  Both exceptional
cases are excluded, proving injectivity.
\end{proof}
